\documentclass[11pt,a4paper]{article}
\usepackage[margin=1in]{geometry}
\usepackage[T1]{fontenc}
\usepackage[utf8]{inputenc}
\usepackage{lmodern}
\usepackage{microtype}
\usepackage{booktabs}
\usepackage{longtable}
\usepackage{hyperref}
\usepackage{parskip}
\usepackage{amsmath}
\title{\textbf{Plants vs. Zombies Deep Learning}\\\large Phase 4 Research Report}
\author{b3d012}
\date{4 September 2026}
\begin{document}
\maketitle

\begin{abstract}
This report documents the offline machine-learning layer and the first bounded
real-game Phase 4 pilots. Phases 1--3.5 are provided by
\texttt{PvZ-AI-Harness} v0.2.3. Phase 4 selects a masked actor--critic baseline,
implements configuration, training, manifests, evaluation, tuning, reward
audit, monitoring, and a live factory, and validates them against a
contract-shaped synthetic environment. Bounded real pilots are reported as
systems evidence only; no learned strategy or convergence claim is made.
\end{abstract}

\section{Relationship to the PvZ AI Harness}
The upstream harness targets Plants vs. Zombies GOTY 1.2.0.1073 on Windows and provides read-only GameState v1, Observation v1, Action v1 with deterministic masking, Controller v1, Environment v1, Reward v1, transition schema v2, baseline policies, and runtime/session safety. The engineering history and validation of those components are documented in the harness repository's cumulative technical development report.

Phase 4 consumes that release as a dependency rather than copying its implementation. The dependency is pinned to \texttt{PvZ-AI-Harness v0.2.3}. A local compatibility gate verifies the expected environment contract before experiments proceed.

\begin{center}
\begin{tabular}{ll}
\toprule
Contract & Initial Phase 4 value \\
\midrule
Harness release & v0.2.3 \\
Observation & schema v1, shape $(5534,)$ \\
Action & schema v1, 541 actions \\
Environment & schema v1 \\
Reward & schema v1 \\
Transition log & schema v2 \\
\bottomrule
\end{tabular}
\end{center}

\section{Phase 4 Objective}
The objective is to determine whether a learned policy can develop useful Plants vs. Zombies strategy from the frozen structured observation/action interface and outperform transparent random-valid and scripted baselines under a reproducible protocol. The model should learn strategy rather than deterministic placement legality or Windows-control mechanics already solved by the harness.

\section{Experimental Principles}
\begin{itemize}
\item Keep the harness fixed within an experiment series.
\item Separate technical runtime interruptions from game outcomes.
\item Record exact Git commits, harness release, schema versions, seeds, hyperparameters, model architecture, timing, and evaluation conditions.
\item Prefer multiple seeds and uncertainty reporting over a single best run.
\item Keep training, evaluation, and hyperparameter selection logically separate.
\item Do not introduce reward terms solely because a metric is convenient to optimize.
\item Preserve raw run manifests and immutable checkpoint provenance outside the source repository when artifacts become large.
\end{itemize}

\section{Phase 4.1: RL Interface and Protocol}
Gymnasium 1.2.2 is used as a narrow adapter contract and sb3-contrib 2.9.0 MaskablePPO as the first algorithm. MaskablePPO natively applies a changing boolean mask during optimization and explicit masked prediction. Plain PPO and A2C lack native masking in the selected stack; DQN may reuse expensive experience better but requires a carefully tested masked target/policy integration; a custom actor--critic creates unnecessary first-baseline risk. Imitation remains a future bootstrap option.

The adapter passes the fixed float vector and semantic action index unchanged. Natural win/loss maps to Gymnasium termination. A configured horizon and every technical interruption map to truncation, with technical provenance retained in \texttt{info}. Thus focus loss is never labeled game loss.

\section{Policy Architecture and Extensibility}
The initial policy/value network uses shared flat Observation v1 input and two hidden layers of 128 units. Tracked 256-by-256 and 512-by-256-by-128 alternatives support controlled comparisons. Structured encoders are deferred because unexplained hard-coded observation slices would be fragile. Algorithm backends expose build, learn, predict, save, and load operations through a registry, isolating algorithms from environment, manifest, and evaluation code.

\section{Training, Checkpoints, and Provenance}
Runs have explicit timestep, episode, and wall-clock limits. Periodic checkpoints and a final checkpoint are written even after a graceful interrupt. Each new directory contains an immutable JSON manifest and resolved YAML plus metrics, TensorBoard, evaluation, and transition subdirectories. Provenance includes both Git identities, all environment schemas, model and algorithm versions, full hyperparameters, five software seed roles, uncontrolled game RNG, device/CUDA information, timing, budget, evaluation protocol, and optional parent lineage. Compatibility is checked before resume; latest checkpoints are never selected silently. A separate immutable \texttt{run\_completion.json} records the final checkpoint path and SHA-256, final step, stop reason, finish time, and optional evaluation reference without rewriting the initial manifest.

\section{Reward and Reward-Hacking Analysis}
The baseline retains harness Reward v1: terminal outcomes dominate, wave progress supplies small shaping, and technical diagnostics carry small penalties. Successful planting and waiting receive no activity bonus. Component logs are aggregated by episode and action; dominance warnings identify returns driven almost entirely by one component. Wave progress, natural win rate, technical failure rate, action balance, and wall time remain external metrics so increasing reward with worsening task performance is visible.

\section{Evaluation and Tuning}
Evaluation is a separate module that supplies masks for every prediction and records all episodes. It reports mean, median, standard deviation, extrema, a normal-approximation 95\% confidence interval, wave progress, and technical truncations. Baselines and learned agents must share level, reward, timing, harness, reset, and horizon. Sequential local Optuna studies are supported because only one client normally exists. While wins are sparse, tuning ranks held-out normalized wave progress then return; final claims require fresh multi-seed confirmation.

\section{Dashboard and Operations}

\section{First Real Phase 4 Pilot}
The first bounded live milestone used the normal Adventure 1--7 daytime
condition, six seed packets, active rows 0--4, a 250 ms decision interval, and
PvZ-AI-Harness v0.2.3. Random-valid completed three horizon-bounded episodes
with returns 0.07, 0.07, and 0.05 and waves 8, 7, and 6. The scripted heuristic
completed three episodes with returns 0.04, 0.0475, and 0.08 and wave progress
9, 9, and 9. Neither baseline had a technical truncation or natural win/loss.

A CPU MaskablePPO MLP-small pilot completed 512 timesteps and wrote checkpoints
at 256 and 512 steps. A separate continuation loaded the final checkpoint and
advanced the model to 1024 total timesteps without overwriting the parent run.
Three live evaluations of the resumed checkpoint returned 0.065, 0.05, and
0.09, reaching waves 9, 7, and 9 with no technical truncation. These results
validate the systems path, not convergence or strategic competence.

A 512-transition live reward audit reconciled episode reward exactly with its
components: wave progress contributed 0.03 and terminal, rejected-action,
controller-failure, plant-not-observed, and state-unavailable components were
zero. The wave-progress dominance warning is expected for a horizon-bounded
episode without a terminal event. Measured throughput was approximately
9--10 thousand strategic steps per hour, implying roughly 10--11 hours for
100k steps before accounting for setup and failures. The game RNG was not
controlled, and three episodes are not a statistically conclusive comparison.

\section{Limitations and Future Work}
The immediate next experiment is a long, multi-seed campaign with fresh held-out
evaluation episodes. It must preserve the immutable harness release, checkpoint
lineage, technical-truncation distinction, and uncontrolled-RNG limitation.
The bounded pilot does not justify claims that the agent has learned competent
Plants vs. Zombies strategy.

\section{Dashboard and Operations}
The Tk dashboard provides Runtime, Board, Agent, Training, Evaluation, Tuning,
and Runs tabs and can inspect immutable run data. Structured view models map
public runtime snapshots, board entities, decoded masked actions, training
history, and evaluation comparisons. It does not duplicate Win32 control or
own configuration truth. The doctor command is read-only. Live training
requires explicit authorization and fails closed at the immutable-release gate.

\section{Planned Research Sequence}
\begin{enumerate}
\item Phase 4.1: RL-facing adapter and experimental protocol.
\item Phase 4.2: first reproducible learned baseline.
\item Phase 4.3: checkpointing, run manifests, recovery, and training metrics.
\item Phase 4.4: algorithm-specific hyperparameter tuning.
\item Phase 4.5: evaluation, ablations, and comparison with frozen baselines.
\end{enumerate}

\section{Reproducibility Record}
Every durable run should record the Phase 4 Git SHA, resolved harness release/SHA, complete environment contract, algorithm/library versions, model architecture, hyperparameters, random seeds, level/episode configuration, reward specification, timing, training budget, hardware/device information, checkpoint identity, and evaluation results.

\section{Known Initial Constraints}
The harness currently requires a manually prepared level start, explicit active rows for levels with inactive lanes, and an authoritative injected terminal detector for natural win/loss. The project therefore must distinguish true environment outcomes from configured truncations and technical runtime interruptions. Automated menu navigation and Adventure progression are outside the initial Phase 4 scope.

\section{Validation and Results}
Offline tests cover contract compatibility, adapter spaces/masks and outcome mapping, configuration rejection, manifests and immutability, completion checkpoint hashing, resume mismatch, evaluation aggregation, reward totals, tuning search-space validity, dashboard models, release-gated live construction, recovery decisions, SB3 metric extraction, and MaskablePPO build/learn/save/load/masked prediction. The bounded real pilots are recorded as PILOT evidence in the results ledger; they validate lifecycle, masked actions, checkpointing, resume, evaluation, reward accounting, and throughput, but do not establish statistical performance or learned competence.

\section{Threats to Validity and Next Work}
Game RNG is uncontrolled; the pilot sample is tiny; on-policy sample cost may be prohibitive; and Reward v1 may be too sparse. Adventure 1-5 remains excluded because it is Wall-nut Bowling rather than a regular economy level. The v0.2.3 harness supplies validated reset and managed pickup services, while Shovel and speed remain deferred. Consequently the implementation is operational for bounded Phase 4 pilots, but the recommended long campaign and statistically meaningful comparison remain future work.

\end{document}
